(1) 自由电荷Gauss定律给 $D_r(2\pi rL)=Q$，故 $D_r=Q/(2\pi Lr)$，$E_r=Q a^s/(2\pi L\varepsilon_a r^{s+1})$。介质内部 $\nabla\cdot\mathbf D=0$，而 $\nabla\cdot\mathbf E$ 可不为零，反映非均匀极化束缚电荷。

(2) 定义 $I_s=\int_a^b a^sr^{-s-1}dr$，对 $s\ne0$ 得 $I_s=[1-(a/b)^s]/s$，对 $s=0$ 得 $I_0=\ln(b/a)$。因此 $V=QI_s/(2\pi L\varepsilon_a)$，$\boxed{C=2\pi L\varepsilon_a/I_s}$；$s\to0$ 时由指数导数恢复对数型普通同轴电容。

(3) 线性电介质的静电场能 $U=\tfrac12\int\mathbf E\cdot\mathbf D\,d^3x=\tfrac12\int_a^b E_rD_r(2\pi Lr)dr=Q^2I_s/(4\pi L\varepsilon_a)=Q^2/(2C)$。固定几何和介电分布从零充到 $Q$ 的外电源功 $\int_0^Q V(q)dq$ 也是同一结果。

(4) 局部场比为 $E(a)/E(b)=(b/a)^{s+1}$。$I_s=\int_a^b (a/r)^s dr/r$ 对 $s$ 单调减小，故 $C$ 随 $s$ 单调增大。$s\to+\infty$ 时 $(a/b)^s\to0$，$I_s\sim1/s$，所以 $C\sim2\pi L\varepsilon_a s$；该数学极限需要材料仍满足给定线性介电律。
